\section*{Contenuti formativi previsti}
% Personalizzare indicando le tecnologie e gli ambiti di interesse dello stage
Durante questo progetto di stage lo studente avrà occasione di approfondire le proprie conoscenze nell'ambito dello sviluppo di applicazioni web full-stack moderne, dell'architettura di piattaforme SaaS e dell'infrastruttura cloud che le supporta. In particolare, gli ambiti formativi previsti sono i seguenti:

\begin{itemize}
    \item \textbf{Sviluppo backend con Node.js 24 LTS e NestJS 11}: progettazione e implementazione di API REST seguendo il paradigma MVC, definizione di controller, service e module, gestione delle dipendenze e testing delle API.

    \item \textbf{Persistenza dati con PostgreSQL 18}: modellazione di entità, scrittura di query, gestione delle migrazioni e integrazione con il backend tramite Supabase.

    \item \textbf{Autenticazione e autorizzazione}: implementazione di flussi di autenticazione basati su JWT, gestione di utenti e ruoli, utilizzo di middleware, guard e validatori.

    \item \textbf{Sviluppo frontend con Next.js 16 e Tailwind CSS 4}: realizzazione di interfacce utente, gestione del routing, fetching dei dati, gestione dello stato applicativo, stilizzazione utility-first e integrazione con il backend.

    \item \textbf{Architetture SaaS e multitenancy}: comprensione dei principi alla base delle piattaforme SaaS, strategie di separazione dei dati e della logica applicativa per la gestione di più clienti.

    \item \textbf{Containerizzazione con Docker Engine 29}: utilizzo di un ambiente di sviluppo Docker pre-configurato, comprensione di Docker Compose, networking tra container e gestione dei volumi.

    \item \textbf{Orchestrazione e cloud}: introduzione a Kubernetes 1.36 nelle distribuzioni k3s e k8s tramite l'ambiente di produzione AWS pre-impostato, comprensione del flusso di deploy dei servizi backend e dell'utilizzo di servizi cloud quali S3 per la gestione dei file; introduzione a Vercel come piattaforma di deploy dei frontend Next.js, con preview environments e flusso di deploy continuo.

    \item \textbf{Architetture scalabili}: introduzione alla messaggistica asincrona con RabbitMQ 4, gestione di processi in background, principi di microservizi e utilizzo di un API Gateway (Kong Gateway 3.10 LTS).

    \item \textbf{Caching e ottimizzazione delle performance}: utilizzo di Redis 8 per il caching applicativo e analisi delle strategie di ottimizzazione lato backend.

    \item \textbf{Pratiche di sviluppo professionale}: utilizzo di sistemi di versionamento (Git), redazione di documentazione tecnica e attività di code review nell'ambito di un team di sviluppo.
\end{itemize}
\newpage